\documentclass{scrreprt}
\usepackage{listings}
\usepackage{underscore}
\usepackage{graphicx}
\usepackage[bookmarks=true]{hyperref}
\usepackage[utf8]{inputenc}
\usepackage[english]{babel}
\usepackage{hyperref}
\usepackage{tabularx}
\usepackage{imakeidx}
\makeindex[intoc]
\usepackage{graphicx}
\graphicspath{ {./Pictures/} }

\hypersetup{
    bookmarks=false,    % show bookmarks bar?
    pdftitle={Software Requirement Specification},    % title
    pdfauthor={Sara Lindholm},                     % author
    pdfsubject={TeX and LaTeX},                        % subject of the document
    pdfkeywords={TeX, LaTeX, graphics, images}, % list of keywords
    colorlinks=true,       % false: boxed links; true: colored links
    linkcolor=blue,       % color of internal links
    citecolor=black,       % color of links to bibliography
    filecolor=black,        % color of file links
    urlcolor=purple,        % color of external links
    linktoc=page            % only page is linked
}%
\def\myversion{1.2}
\title{Software Requirements Specification}
\author{Sara Lindholm \\ Patrick Asman \\ Jacob Rhan}

\date{October 2020}

\newcolumntype{s}{>{\hsize=3.75\hsize}X}
\newcolumntype{m}{>{\hsize=7\hsize}X}
\newcolumntype{b}{>{\hsize=11\hsize}X}

\begin{document}

\begin{titlepage}
    \begin{center}
    \begin{bfseries}
        \Huge{SOFTWARE REQUIREMENTS SPECIFICATION}\\
        \vspace{1.5cm}
        \LARGE Patrick Asman \\ Sara Lindholm \\ Jacob Rhan \\
        \vspace{1.5cm}
        %Submitted to : - \\Lecturer\\
        %\vspace{1.5cm}
        \today\\
        %\vspace{5cm}
        \vspace{1.5cm}
        {Version \myversion}\\
        \vfill
        \includegraphics[width=\linewidth]{Pictures/logo.png} \\
    \end{bfseries}
    \end{center}
\end{titlepage}

\tableofcontents

\chapter*{Revision History}
\begin{center}
\begin{tabular}{|p{1,9cm}|p{1,5cm}|p{6cm}|p{4cm}|}
 \hline
 \textbf{Date} & \textbf{Version} & \textbf{Description} & \textbf{Author} \\ 
 \hline
 2020-10-xx & 0.9 & Initial SRS structure & Sara Lindholm \\
 \hline
 2020-10-15 & 1.0 & First Draft & Patrick Asman, Sara Lindholm, Jacob Rhan \\ 
 \hline
 2020-11-11 & 1.1 & Updated and added functional requirements & Sara Lindholm \\
 \hline
 2020-11-17 & 1.2 & Added category descriptions & Jacob Rhan, Sara Lindholm\\
 \hline
\end{tabular}
\end{center}



\chapter{Introduction}
This section gives an overview of everything included in this document. The purpose and scope for this document is described and a list of abbreviations, acronyms and definitions is provided. 

\section{Purpose}
The purpose of this document is to give a detailed description of \textcolor{red}{[system name?]}. It will illustrate the purpose and features of the system. It will also explain system constraints, interfaces of the system and interactions with external applications. This document is intended primarily to be proposed to the stakeholders for approval and as a reference for the developers.

\section{Scope}
\textcolor{red}{[System name?]} is a web-based application to aid Region Östergötland's health care staff in overseeing and handling many patients. The system will make it easier for the health care staff to oversee many patients' measured health data and prioritize which patients are in most need of care. The system will have customizable overviews, customizable alert-levels and sorting and filtering functions to make prioritization easier. The system will in addition to a better overview of all patients be able to give more detailed overview about the individual patients' health information, such as current medications, health conditions and the measured data. 


\section{Definitions, acronyms and abbreviations}
\begin{center}
\begin{tabularx}{\linewidth}{| l | X |}
 \hline
 \textbf{Term} & \textbf{Definition} \\ 
 \hline
  User  & Health care staff who interacts with the web application.  \\
  \hline
  Admin/Administrator & System administrator who is given specific permission for managing and controlling the system  \\ 
  \hline
  RÖ  & Region Östergötland  \\ 
  \hline
\end{tabularx}
\end{center}

\section{References}
IEEE Recommended Practice for Software Requirements Specifications," in IEEE Std 830-1998 , vol., no., pp.1-40, 20 Oct. 1998, doi: 10.1109/IEEESTD.1998.88286.

\section{Overview}
The second chapter of the document, Overall Description, gives an overview of the functionality of the product.


The third chapter, Requirements Specification, gives a more  detailed description of the system and its features. 

\chapter{Overall Description}
This section will present an overview of the entire system. The system's basic functionality and how it interacts with other systems will be introduced. % The constraints and assumptions for the system will also be presented. 

\section{Product Perspective}
\textbf{TODO:} fix. \\ put the product into perspective with other related products
\begin{itemize}
    \item web application
    \item communicates with openEHR (store and recieve health data?), OAuth, RÖ's personell lists, Swedish population register and online health care services (such as 1177, Fass etc).
\end{itemize}


\section{Product functions}
This web application will aid health care staff in handling many patients and their health data. The system will have an overview in which the user will be able to see all patients in their care unit. These patients can then be sorted and filtered by different criteria, which will aid the user in prioritizing which patients are in most need of care. 


The user will also be able to show more detailed information about the individual patients should they choose to click on a patient. In the individual patient view the user will get access to more information than what was available in the overview of all patients. The individual patient may have measured health data, medical information and history, and scheduled appointments which will be displayed in different tabs on the web application.


The application will also be able to send out notifications to the users when there are updates on patients. The specifics for notifications are to be set by the admin so as to avoid being overwhelmed by notifications. 

\section{User Characteristics}
There are two intended users of the system: users and administrators. These two types of users will have different roles and therefore not have access to all the same functionalities. Some of the customizable elements shall only be customizable by an admin and the admin will also be the one to set the criteria for the different alert levels as well as the notifications.


\section{Constraints}


\section{Assumptions and dependencies}
%"This subsection of the SRS should list each of the factors that affect the requirements stated in the SRS. These factors are not design constraints on the software but are, rather, any changes to them that can affect the requirements in the SRS. For example, an assumption may be that a specific operating system will be available on the hardware designated for the software product. If, in fact, the operating system is not available, the SRS would then have to change accordingly."

\section{Apporting of requirements}
%There are some requirements that could be delayed until future versions of the systems if the project is delayed.

\chapter{Specific requirements}
This chapter contains the software requirements and gives a detailed
description of the system and its features. 

\section{External interface requirements}
\textbf{TODO: fix}
\begin{itemize}
    \item openEHR
    \item OAuth
    \item Region Östergötland's personell lists 
    \item Swedish population register 
    \item online health care services (e.g. rådgivningsstöd, 1177, vårdhandboken, FASS, internetmedicin) 
\end{itemize}

%\subsection{User interfaces}
%\subsection{Software interfaces}
%\subsection{Communication interfaces}

\section{Functional requirements}
The following section includes the requirements that specify all the fundamental actions of the software system. The section is divided into four main categories: "Admin", "Patient", "System" and "Översikt". These categories are then divided into subcategories based on the different views in each category. 
\\ \\
To clarify how the requirements presented in tables see explanation below. 

\begin{center}
\begin{tabularx}{\linewidth}{| s | m | m | b | s |}
\hline
\textbf{ID:} & \textbf{Title} & \textbf{Precondition} & \textbf{Description} & \textbf{Priority} \\ 
\hline
FRXXX & 
Title to briefly explain the requirement & 
Explains any preconditions for requirement & 
Description of what will happen in the system.  
    \newline \textbf{Input:} What triggers the function  
    \newline \textbf{Action:} How the system responds to the input 
    \newline \textbf{Output:} The function's output &
Priority level (Low-Medium-High) \\ 
\hline
\end{tabularx}
\end{center}

\subsection{Admin}
In section 3.2.1. we describe functional requirements for the Admin. A precondition is that the admin is logged in to the system. 
\subsubsection{Översikt/Patienter}
A precondition is that the admin is logged in and on the view “Oversikt/Patienter”
\begin{center}
\begin{tabularx}{\linewidth}{| l | X | X | l |}
\hline
\textbf{ID:} & \textbf{Statement} & \textbf{Description} & \textbf{Priority} \\ 
\hline
FR001 & 
The admin shall be able to customize and save views of the system (e.g. diabetes over 30 years, a view that filter out all diabetes patients that are older than 30 years) & 
\textbf{Input:} Selected filter settings by Admin. \newline \textbf{Action:} System generating new view. \newline \textbf{Output:} New view for users. &
Medium \\ 
\hline
FR002 & The "Översiktsinställningar" shall only be customizable by the administrator. & \textbf{Input:} Selected settings by Admin. \newline \textbf{Action:} System generating new views. \newline \textbf{Output:} Updated views for users. & Medium \\
\hline
FR003 & 
The admin shall be able to determine when notifications shall be sent to users. & \textbf{Input:} Selecting the notifications to receive for a specific patient (e.g. Want to receive notifications when patient goes from one priority level to another one).
\newline \textbf{Action:} System sets notifications for patient.
\newline \textbf{Output:} "Notification prefences updated" message. & 
Medium \\ 
\hline
\end{tabularx}
\end{center}

\subsection{Patient}
In section 3.2.2. we describe functional requirements for the User. A precondition is that the user is logged in to the system and, on the view “Patient”
\subsubsection{Admin}
A precondition is that the admin is logged in and on the view “Patient/Admin” 
\begin{center}
\begin{tabularx}{\linewidth}{| l | X | X | l |}
\hline
\textbf{ID:} & \textbf{Statement} & \textbf{Description} & \textbf{Priority} \\ 
\hline
FR004 & The system shall log users accessing patient information. & \textbf{Input:} User opens individual patient view / User ID.
\newline \textbf{Action:} System logs ID and updates activity log.\newline \textbf{Output:} Updated activity log. & High \\ 
\hline
FR005 & Admin shall configure at what level (based on health care status, health data etc) a measurement shall be to trigger a priority level change about the patient & \textbf{Input:} Limit values.
\newline \textbf{Action:} Save patient specific prioritization limits.
\newline \textbf{Output:} Settings displayed. &  High \\ 
\hline
FR006 & The system shall log users updating patient information. & \textbf{Input:} User updates patient information (e.g measurements).
\newline \textbf{Action:} System logs ID and updates activity log.
\newline \textbf{Output:} Updated activity log. & High \\ 
\hline
\end{tabularx}
\end{center}

\subsubsection{Alla flikar}
\begin{center}
\begin{tabularx}{\linewidth}{| l | X | X | l |}
\hline
\textbf{ID:} & \textbf{Statement} & \textbf{Description} & \textbf{Priority} \\ 
\hline
FR007 & In the view of the patient the current diagnosis/diagnoses shall be presented in the heading  & \textbf{Input:} Click in on individual patient.
\newline \textbf{Action:} Data loading.
\newline \textbf{Output:} Diagnosis information displayed. & Medium \\ 
\hline
FR008 & The user shall be able to switch to Overview view. & \textbf{Input:} User clicks on the Home button.
\newline \textbf{Action:} Loading Overview view.
\newline \textbf{Output:} System displays Overview View. & High \\ 
\hline
\end{tabularx}
\end{center}

\subsubsection{Kalender}
A precondition is that the user is logged in and on the view “Patient” and, on the tab “Kalender”
\begin{center}
\begin{tabularx}{\linewidth}{| l | X | X | l |}
\hline
\textbf{ID:} & \textbf{Statement} & \textbf{Description} & \textbf{Priority} \\ 
\hline
FR009 & In the patient view under the calendar tab, the calendar activities must be color-coded based on the activity type. & \textbf{Input:} input.
\newline \textbf{Action:} action.
\newline \textbf{Output:} output. & Low \\ 
\hline
FR010 & In the patient view under the calendar tab, the Calendar must be able to be linked to "timebooks" / “tidböcker”. & \textbf{Input:} input.
\newline \textbf{Action:} action.
\newline \textbf{Output:} output. & Low \\ 
\hline
FR011 & In the calendar view of a patient the user shall for an upcoming activity be presented with the date, time, description, place, and comments. & \textbf{Input:} Click tab "Patient/Kalender".
\newline \textbf{Action:} Load activity information.
\newline \textbf{Output:} Display upcoming  activity information. & Low \\ 
\hline
FR012 & The user shall be able to add activities to a patient's calendar.  & \textbf{Input:} input.
\newline \textbf{Action:} action.
\newline \textbf{Output:} output. & Low \\ 
\hline
\end{tabularx}
\end{center}

\subsubsection{Mätvärden}
A precondition is that the user is logged in and on the view “Patient” and, on the tab “Mätvärden”
\begin{center}
\begin{tabularx}{\linewidth}{| l | X | X | l |}
\hline
\textbf{ID:} & \textbf{Statement} & \textbf{Description} & \textbf{Priority} \\ 
\hline
FR013 & In the view of the individual patient's measured data the user of the system shall be able to zoom in and out in the view.  & \textbf{Input:} Press zoom in/out button.
\newline \textbf{Action:} Resizing diagram and list of measurements.
\newline \textbf{Output:} Resized diagrams and list of measurements. & Low \\ 
\hline
FR014 & It must be specified during what period the shown measurments have been collected. & \textbf{Input:} Click tab "Patient/Mätvärden".
\newline \textbf{Action:} Measurment loading.
\newline \textbf{Output:} Diagrams of measurments with dates on axis. & Low \\ 
\hline
FR015 & The system shall present measurement history in a diagram. & \textbf{Input:} Click tab "Patient/Mätvärden".
\newline \textbf{Action:} Data loading.
\newline \textbf{Output:} Information displayed on view. & Medium \\ 
\hline
FR016 & The user shall be able to accept or deny unexpected measured values submitted by a patient.  & \textbf{Input:} Click on unexpected measurement value.
\newline \textbf{Action:} System accepts or denies measured value.
\newline \textbf{Output:} Displays changes. & Medium \\ 
\hline
FR017 & The user shall be able to update a patients measured values.  & \textbf{Input:} Clicks on measurement to update.
\newline \textbf{Action:} System updates measured value.
\newline \textbf{Output:} Displays value. & High \\ 
\hline
\end{tabularx}
\end{center}

\subsubsection{Översikt}
A precondition is that the user is logged in and on the view “Patient” and, on the tab “Översikt”
\begin{center}
\begin{tabularx}{\linewidth}{| l | X | X | l |}
\hline
\textbf{ID:} & \textbf{Statement} & \textbf{Description} & \textbf{Priority} \\ 
\hline
FR018 & Information about allergies and other care restrictions must be visible in patient overview. & \textbf{Input:} Click tab "Patient/Översikt".
\newline \textbf{Action:} Data loading.
\newline \textbf{Output:} Information displayed on view. & High \\ 
\hline
FR019 & Personal information (gender, age, personal number) of the patient shall be included in the view. & \textbf{Input:} Click tab "Patient/Översikt".
\newline \textbf{Action:} Data loading.
\newline \textbf{Output:} Information displayed. & Medium \\ 
\hline
FR020 & Contact information and emergency contacts, if available, shall be included in the view. & \textbf{Input:} Click tab "Patient/Översikt".
\newline \textbf{Action:} Data loading.
\newline \textbf{Output:} Information displayed. & Medium \\ 
\hline
FR021 & At the top of the view, the patient's current diagnosis/diagnoses shall be visible together with a collapsible list of past diagnosis/noses & \textbf{Input:} Click tab "Patient/Översikt".
\newline \textbf{Action:} Data loading.
\newline \textbf{Output:} Information displayed at the top of the view. & Medium \\ 
\hline
FR022 & In the overview of a patient the user shall in the middle of the view be presented with measurements and a seven-day measurement average & \textbf{Input:} Click tab "Patient/Översikt".
\newline \textbf{Action:} Data loading.
\newline \textbf{Output:} Display measurements and a seven day average & Medium \\ 
\hline
FR023 & In the overview of a patient the user shall at the top of the view be presented with the patient’s current medication(s) & \textbf{Input:} Click tab "Patient/Översikt".
\newline \textbf{Action:} Data loading.
\newline \textbf{Output:} Information displayed at the top of the view. & Medium \\ 
\hline
FR024 & In the overview of a patient the user shall at the bottom of the view be presented with responsible healthcare provider(s) and associated hospital or clinic.  & \textbf{Input:} Click tab "Patient/Översikt".
\newline \textbf{Action:} Data loading.
\newline \textbf{Output:} Information displayed at the bottom of the view.  & Low \\
\hline
\end{tabularx}
\begin{tabularx}{\linewidth}{| l | X | X | l |}
\hline
FR025 & The user shall, if necessary, be presented with "Uppmärksamhetssymbolen" in the overview of a patient & \textbf{Input:} Click tab "Patient/Översikt".
\newline \textbf{Action:} Data loading.
\newline \textbf{Output:} Information displayed in the view. & Medium \\ 
\hline
\end{tabularx}
\end{center}

\subsection{System}
In section 3.2.3. we describe functional requirements for the system. A precondition is that the user is logged in to the system.
\begin{center}
\begin{tabularx}{\linewidth}{| l | X | X | l |}
\hline
\textbf{ID:} & \textbf{Statement} & \textbf{Description} & \textbf{Priority} \\ 
\hline
FR026 & The system shall notify the user if patients have not submitted medical data according to their individual health care plan. & \textbf{Input:} Time limit, according to individual health care plan, expires.
\newline \textbf{Action:} Notification printing and sending to responsible user.
\newline \textbf{Output:} The system sends notification to user. & High \\ 
\hline
FR027 & The system shall provide the user with instructions on how to use the system. For example, if a field is to be filled in, there is information about what is to be specified in the field. & \textbf{Input:} User clicks on button "Help".
\newline \textbf{Action:} Loading help information message.
\newline \textbf{Output:} The system displays help information message. & Low \\ 
\hline
FR028 & The system shall send a notification when a patient measurement exceeds limit value. & \textbf{Input:} Patient measurement.
\newline \textbf{Action:} Comparing to limit value.
\newline \textbf{Output:} Send notification. & High \\
\hline
FR029 & The user of the system shall be presented with the option to access communication information (i.e. chat, telephone number, contact information) when notified about a patient & \textbf{Input:} Notification sent with patient identification.
\newline \textbf{Action:} Load contact information about patient.
\newline \textbf{Output:} Display patient's contact information in notification if uses so chooses. & Low \\
\hline
\end{tabularx}
\end{center}

\subsection{Översikt}
In section 3.2.4. we describe functional requirements for the User. A precondition is that the user is logged in to the system and, on the view “Översikt”
\subsubsection{Hem}
\begin{center}
\begin{tabularx}{\linewidth}{| l | X | X | l |}
\hline
\textbf{ID:} & \textbf{Statement} & \textbf{Description} & \textbf{Priority} \\ 
\hline
FR030 & When the filter view is open, the user shall be able to filter patients by [age, priority-level, measurements, health care unit, team, disease and patients’conditions among others].  & \textbf{Input:} Filter values from categories [age, priority-level, measurements, health care unit, team, disease and patients’conditions among others].
\newline \textbf{Action:} The system searches for patients matching filter criteria.
\newline \textbf{Output:} The system displays all patients who matched the filters. & High \\ 
\hline
FR031 & The user of the system shall be able to sort patients by [age, priority-level, measurements, health care unit, team, patients’ conditions, disease type, latest date, name, updated by ] & \textbf{Input:} Sorting values from categories [age, priority-level, measurements, health care unit, team, disease, patients’conditions, disease type, latest date, name, updated by].
\newline \textbf{Action:} The system rearranges the list of patients according to sorting criteria.
\newline \textbf{Output:} The system displays all patients according to the sorting criteria. & High \\ 
\hline
FR032 & When the user tries to access patient data from other departments the system shall display a message & \textbf{Input:} Clicks on patient not in user's department.
\newline \textbf{Action:} The system confirms that user wants to access patient data not in their department.
\newline \textbf{Output:} Display message about accessing patient data outside user's department. & High \\ 
\hline
FR033 & The overview page shall be shown to users as starting page after login. & \textbf{Input:} User logs in.
\newline \textbf{Action:} Loading starting page (Overview view).
\newline \textbf{Output:} System displays Overview view. & High \\ 
\hline
\end{tabularx}
\begin{tabularx}{\linewidth}{| l | X | X | l |}
\hline
FR034 & It must be visible which parameter the sorting is based on. & \textbf{Input:} Click sort.
\newline \textbf{Action:} Data loading.
\newline \textbf{Output:} Information displayed in relation to search result. & Low \\ 
\hline
FR035 & Overview-dashboard shall display information according to the admin preference settings (e.g name, personal number, prioritization, diagnosis, health care provider responsible for patient, last updated, updated by) & \textbf{Input:} Selecting which "Översiktsinställningar" to show in overview.
\newline \textbf{Action:} The system updates the settings.
\newline \textbf{Output:} The system displays all the selected settings for the patients. & Medium \\ 
\hline
FR036 & The user shall be given a choice between accessing the data or not before getting access to said data. & \textbf{Input:} Clicking on a patient in overview.
\newline \textbf{Action:} The system checks if user wants to access patient data.
\newline \textbf{Output:} Display message checking if user wants to access data. & High \\ 
\hline
FR037 & The user shall only get access to the data if the user chooses to access it. & \textbf{Input:} Clicking on "Bekräfta" button to access patient data.
\newline \textbf{Action:} The system confirms that data shall be displayed.
\newline \textbf{Output:} The system displays patient data. & High \\ 
\hline
FR038 & When search filtering, the active list shall be filtered dynamically as the user types in the search query. & \textbf{Input:} Text string in diagnosis filter.
\newline \textbf{Action:} System checks for diagnosis that matches text string.
\newline \textbf{Output:} Displays diagnosis that matches string. & Low \\
\hline
FR039 & The active hover state of a notification-symbol shall present relevant text information (e.g. Unexpected measurement value, or delayed time) for the user   & \textbf{Input:} Mouse hover over notification.
\newline \textbf{Action:} System loads data.
\newline \textbf{Output:} Displays relevant infromation. & Medium \\ 
\hline
FR040 & The user shall be able to switch to patient view. & \textbf{Input:} User clicks on a patient.
\newline \textbf{Action:} Loading Patient view.
\newline \textbf{Output:} System displays Patient View. & High \\ 
\hline
\end{tabularx}
\begin{tabularx}{\linewidth}{| l | X | X | l |}
\hline
FR041 & The system shall have a symbol that indicates what type of notification it is [1. Exceeded limit value, 2. Delayed measurment] & \textbf{Input:} Click on "Översikt/Notislogg".
\newline \textbf{Action:} Load symbols.
\newline \textbf{Output:} Display symbols. & Medium \\ 
\hline
\end{tabularx}
\end{center}

\subsubsection{Kalender}
A precondition is that the user is logged in, on the view “Översikt” and, on the tab “Kalender”
\begin{center}
\begin{tabularx}{\linewidth}{| l | X | X | l |}
\hline
\textbf{ID:} & \textbf{Statement} & \textbf{Description} & \textbf{Priority} \\ 
\hline
FR042 & In the calendar view of a patient the user shall for an upcoming activity be presented with the date, time, description, place, and comments & \textbf{Input:} Click tab "Översikt/Kalender".
\newline \textbf{Action:} Load activity information.
\newline \textbf{Output:} Display upcoming  activity information. & Low \\ 
\hline
\end{tabularx}
\end{center}

\subsubsection{Notislogg}
A precondition is that the user is logged in, on the view “Översikt” and, on the tab “Notislog”
\begin{center}
\begin{tabularx}{\linewidth}{| l | X | X | l |}
\hline
\textbf{ID:} & \textbf{Statement} & \textbf{Description} & \textbf{Priority} \\ 
\hline
FR043 & The system shall have a symbol that indicates what type of item it is [1. Exceeded limit value, 2. Delayed measurment] & \textbf{Input:} Click on "Översikt/Notislogg".
\newline \textbf{Action:} Load symbols.
\newline \textbf{Output:} Display symbols. & Medium \\
\hline
FR044 & The user shall be given a choice between accessing the data or not before getting access to said data. & \textbf{Input:} Clicking on a patient in overview.
\newline \textbf{Action:} The system checks if user wants to access patient data.
\newline \textbf{Output:} Display message checking if user wants to access data. & High \\
\hline
FR045 & The user shall only get access to the data if the user chooses to access it. & \textbf{Input:} Clicking on "Bekräfta" button to access patient data.
\newline \textbf{Action:} The system confirms that data shall be displayed.
\newline \textbf{Output:} The system displays patient data. & High \\ 
\hline
FR046 & It must be visible which parameter the sorting is based on. & \textbf{Input:} Click sort.
\newline \textbf{Action:} Data loading.
\newline \textbf{Output:} Information displayed in relation to search result. & Low \\ 
\hline
FR047 & The active hover state of a notification-symbol shall present relevant text information (e.g. Unexpected measurement value, or delayed time) for the user   & \textbf{Input:} Mouse hover over notification.
\newline \textbf{Action:} System loads data.
\newline \textbf{Output:} Displays relevant information. & Medium \\ 
\hline
\end{tabularx}
\begin{tabularx}{\linewidth}{| l | X | X | l |}
\hline
FR048 & The user of the system shall be able to sort patients by [age, priority-level, measurements, health care unit, team, patients’ conditions, disease type, latest date, name, updated by] & \textbf{Input:} Sorting values from categories [age, priority-level, measurements, health care unit, team, disease, patients’conditions, disease type, latest date, name, updated by]. .
\newline \textbf{Action:} The system rearranges the list of patients according to sorting criteria.
\newline \textbf{Output:} The system displays all patients according to the sorting criteria. & High \\ 
\hline
FR049 & When the user tries to access patient data from other departments the system shall display a message & \textbf{Input:} Clicks on patient not in user's department.
\newline \textbf{Action:} The system confirms that user wants to access patient data not in their department.
\newline \textbf{Output:} Display message about accessing patient data outside user's department. & High \\ 
\hline
FR050 & When search filtering, the active list shall be filtered dynamically as the user types in the search query     & \textbf{Input:} Text string in diagnosis filter.
\newline \textbf{Action:} System checks for diagnosis that matches text string.
\newline \textbf{Output:} Displays diagnosis that matches string. & Medium \\ 
\hline
FR051 & The system shall filter patients by [age, alert-level, measurements, health care unit, team, disease and patients’conditions among others].  & \textbf{Input:} Filter values from categories [age, alert-level, measurements, health care unit, team, disease and patients’conditions among others].
\newline \textbf{Action:} The system searches for patients matching filter criteria.
\newline \textbf{Output:} The system displays all patients who matched the filters. & High \\ 
\hline
FR052 & The user shall be able to switch to patient view. & \textbf{Input:} User clicks on a patient.
\newline \textbf{Action:} Loading Patient view.
\newline \textbf{Output:} System displays Patient View. & High \\ 
\hline
\end{tabularx}
\end{center}

\subsubsection{Patienter}
A precondition is that the user is logged in, on the view “Översikt” and, on the tab “Patienter”
\begin{center}
\begin{tabularx}{\linewidth}{| l | X | X | l |}
\hline
\textbf{ID:} & \textbf{Statement} & \textbf{Description} & \textbf{Priority} \\ 
\hline
FR053 & The system shall filter patients by [age, alert-level, measurements, health care unit, team, disease and patients’conditions among others]. The system shall filter patients by [age, alert-level, measurements, health care unit, team, disease and patients’conditions among others]. & \textbf{Input:} Filter values from categories [age, alert-level, measurements, health care unit, team, disease and patients’conditions among others].
\newline \textbf{Action:} The system searches for patients matching filter criteria.
\newline \textbf{Output:} The system displays all patients who matched the filters. & High \\ 
\hline
FR054 & The user shall be given a choice between accessing the data or not before getting access to said data. & \textbf{Input:} Clicking on a patient in overview.
\newline \textbf{Action:} The system checks if user wants to access patient data.
\newline \textbf{Output:} Display message checking if user wants to access data. & High \\ 
\hline
FR055 & The user shall only get access to the data if the user chooses to access it. & \textbf{Input:} Clicking on "Bekräfta" button to access patient data.
\newline \textbf{Action:} The system confirms that data shall be displayed.
\newline \textbf{Output:} The system displays patient data. & High \\ 
\hline
FR056 & It must be visible which parameter the sorting is based on. & \textbf{Input:} Click sort.
\newline \textbf{Action:} Data loading.
\newline \textbf{Output:} Information displayed in relation to search result. & Low \\ 
\hline
FR057 & When search filtering, the active list shall be filtered dynamically as the user types in the search query. & \textbf{Input:} Text string in diagnosis filter.
\newline \textbf{Action:} System checks for diagnosis that matches text string.
\newline \textbf{Output:} Displays diagnosis that matches string. & Medium \\ 
\hline
\end{tabularx}
\begin{tabularx}{\linewidth}{| l | X | X | l |}
\hline
FR058 & The active hover state of a notification-symbol shall present relevant text information (e.g. Unexpected measurement value, or delayed time) for the user   & \textbf{Input:} Mouse hover over notification.
\newline \textbf{Action:} System loads data.
\newline \textbf{Output:} Displays relevant information. & Medium \\ 
\hline
FR059 & The user of the system shall be able to sort patients by [age, priority-level, measurements, health care unit, team, patients’ conditions, disease type, latest date, name, updated by] & \textbf{Input:} Sorting values from categories [age, priority-level, measurements, health care unit, team, disease, patients’conditions, disease type, latest date, name, updated by].
\newline \textbf{Action:} The system rearranges the list of patients according to sorting criteria.
\newline \textbf{Output:} The system displays all patients according to the sorting criteria. & High \\ 
\hline
FR060 & The system shall have a symbol that indicates what type of notification it is [1. Exceeded limit value, 2. Delayed measurement] & \textbf{Input:} Click on "Översikt/Notislogg".
\newline \textbf{Action:} Load symbols.
\newline \textbf{Output:} Display symbols. & Medium \\ 
\hline
FR061 & When the user tries to access patient data from other departments the system shall display a message & \textbf{Input:} Clicks on patient not in user's department.
\newline \textbf{Action:} The system confirms that user wants to access patient data not in their department.
\newline \textbf{Output:} Display message about accessing patient data outside user's department. & High \\ 
\hline
FR062 & Overview-dashboard shall display information according to the admin preference settings (e.g name, personal number, prioritization, diagnosis, health care provider responsible for patient, last updated, updated by) & \textbf{Input:} Input: Selecting which "Översiktsinställningar" to show in overview.
\newline \textbf{Action:} The system updates the settings.
\newline \textbf{Output:}  The system displays all the selected settings for the patients. & Medium \\ 
\hline
FR063 & The user shall be able to switch to patient view. & \textbf{Input:} User clicks on a patient.
\newline \textbf{Action:} Loading Patient view.
\newline \textbf{Output:} System displays Patient View. & High \\ 
\hline
\end{tabularx}
\end{center}

\section{Performance requirements}
The requirements in this section provide specification of the system performance. 

\begin{center}
\begin{tabularx}{\linewidth}{| l | X |}
 \hline
 \textbf{ID:} & NFR005  \\ 
 \hline
 \textbf{Statement:} & The system shall be able to store and handle 10 000 simultaneous active patients.
 \\ 
 \hline
\end{tabularx}
\end{center}


\section{Software system attributes}
The requirements in this section specify the required usability, security, maintainability and portability.
of the software system.

\subsection{Usability}
\begin{center}
\begin{tabularx}{\linewidth}{| l | X |}
 \hline
 \textbf{ID:} & NFR008  \\ 
 \hline
 \textbf{Statement:} & 90\% of the system users shall perceive the system as easy to use and flexible.
 \\ 
 \hline
\end{tabularx}
\end{center}

\begin{center}
\begin{tabularx}{\linewidth}{| l | X |}
 \hline
 \textbf{ID:} & NFR016  \\ 
 \hline
 \textbf{Statement:} & The system shall be easy to navigate (e.g. low Lostness-score, follow the ISO 9241-11 standard).
 \\ 
 \hline
\end{tabularx}
\end{center}

\begin{center}
\begin{tabularx}{\linewidth}{| l | X |}
 \hline
 \textbf{ID:} & NFR017  \\ 
 \hline
 \textbf{Statement:} & The system shall provide good user experience (follow the ISO 9241-11 standard).
 \\ 
 \hline
\end{tabularx}
\end{center}

\begin{center}
\begin{tabularx}{\linewidth}{| l | X |}
 \hline
 \textbf{ID:} & NFR018  \\ 
 \hline
 \textbf{Statement:} & The system shall be accessible for people with disabilities and follow WCAG guidelines. 
 \\ 
 \hline
\end{tabularx}
\end{center}

\begin{center}
\begin{tabularx}{\linewidth}{| l | X |}
 \hline
 \textbf{ID:} & NFR012  \\ 
 \hline
 \textbf{Statement:} & The Google Material Design system themed with Region Östergötland’s color palette shall be used.
 \\ 
 \hline
\end{tabularx}
\end{center}

\subsection{Security}

\begin{center}
\begin{tabularx}{\linewidth}{| l | X |}
 \hline
 \textbf{ID:} & NFR002  \\ 
 \hline
 \textbf{Statement:} & OAuth shall be used for login and authentication.
 \\ 
 \hline
\end{tabularx}
\end{center}

\subsection{Maintainability}

\begin{center}
\begin{tabularx}{\linewidth}{| l | X |}
 \hline
 \textbf{ID:} & NFR013  \\ 
 \hline
 \textbf{Statement:} & The system shall enable use of openEHR platforms as storage backend for EHR content (e.g. measurements, content from forms filled out by the patient and care plan).
 \\ 
 \hline
\end{tabularx}
\end{center}

\begin{center}
\begin{tabularx}{\linewidth}{| l | X |}
 \hline
 \textbf{ID:} & NFR014  \\ 
 \hline
 \textbf{Statement:} & The system shall be integrable with RÖ’s systems for personell lists from their database.
 \\ 
 \hline
\end{tabularx}
\end{center}

\begin{center}
\begin{tabularx}{\linewidth}{| l | X |}
 \hline
 \textbf{ID:} & NFR015  \\ 
 \hline
 \textbf{Statement:} & The system shall be integrable with the Swedish population register (folkbokföringsregister). 
 \\ 
 \hline
\end{tabularx}
\end{center}

\begin{center}
\begin{tabularx}{\linewidth}{| l | X |}
 \hline
 \textbf{ID:} & NFR019  \\ 
 \hline
 \textbf{Statement:} & Online health care services (i.e. rådgivningsstöd, 1177, vårdhandboken, FASS, internetmedicin) shall be integrated in the system.
 \\ 
 \hline
\end{tabularx}
\end{center}

\subsection{Portability}

\begin{center}
\begin{tabularx}{\linewidth}{| l | X |}
 \hline
 \textbf{ID:} & NFR001  \\ 
 \hline
 \textbf{Statement:} & The system shall be OS independent.
 \\ 
 \hline
\end{tabularx}
\end{center}

\begin{center}
\begin{tabularx}{\linewidth}{| l | X |}
 \hline
 \textbf{ID:} & NFR004  \\ 
 \hline
 \textbf{Statement:} & It shall be possible to deploy the system using Docker orchestrated by Kubernetes.
 \\ 
 \hline
\end{tabularx}
\end{center}

\begin{center}
\begin{tabularx}{\linewidth}{| l | X |}
 \hline
 \textbf{ID:} & NFR006  \\ 
 \hline
 \textbf{Statement:} & Standardized terms and national language for special purposes (LSP) shall be used. 
 \\ 
 \hline
\end{tabularx}
\end{center}

\begin{center}
\begin{tabularx}{\linewidth}{| l | X |}
 \hline
 \textbf{ID:} & NFR007  \\ 
 \hline
 \textbf{Statement:} & The system shall be written in TypeScript (instead of JavaScript).
 \\ 
 \hline
\end{tabularx}
\end{center}

\begin{center}
\begin{tabularx}{\linewidth}{| l | X |}
 \hline
 \textbf{ID:} & NFR009  \\ 
 \hline
 \textbf{Statement:} & The system shall implement a scalable UI that works seamlessly across a broad range of screen sizes.
 \\ 
 \hline
\end{tabularx}
\end{center}


\section{Other requirements}

\begin{center}
\begin{tabularx}{\linewidth}{| l | X |}
 \hline
 \textbf{ID:} & NFR003  \\ 
 \hline
 \textbf{Statement:} & The server-side functionality shall be based on open source parts.
 \\ 
 \hline
\end{tabularx}
\end{center}

\begin{center}
\begin{tabularx}{\linewidth}{| l | X |}
 \hline
 \textbf{ID:} & NFR010  \\ 
 \hline
 \textbf{Statement:} & The system shall be developed in open source, Apache 2 license.
 \\ 
 \hline
\end{tabularx}
\end{center}

\begin{center}
\begin{tabularx}{\linewidth}{| l | X |}
 \hline
 \textbf{ID:} & NFR011  \\ 
 \hline
 \textbf{Statement:} & The graphical interface components shall be continuously and openly documented/tested/published through Storybook.
 \\ 
 \hline
\end{tabularx}
\end{center}


\appendix
\chapter{Appendix Title}

%\printindex

\end{document}